\centerline{\bf \Huge Олимпиада I}
\title{Олимпиада I}

1. Существует ли точный квадрат, десятичная запись которого состоит из 993 нулей, 7 единиц и 11 двоек?

2. Очир купил в магазине n аквариумов и поместил у себя дома. В первом аквариуме у него плавала одна рыбка, во втором 3 рыбки, в третьем 5 рыбок и т.д. в $n-$ом акавируме плавала $2n-1$ рыбка. Докажите, что суммарное количество рыбок в аквариумах в $n$ раз больше чем аквариумов.

3. Шестерым Лососям надоело мыть Шиншилл и они решили помыть друг друга, но не все на это согласились. Докажите, что среди них либо найдутся трое, которые помыли друга друга, либо трое которые не помыли друг друга. 

4. Докажите, что $36^{36} + 41^{41}$ делится на $77$.

5. Найдите последнюю цифру чисел $7^{7}$, $7^{7^{7}}$, $7^{7^{7^{7}}}$.

6. Между двумя любыми животноводческими стоянками в Калмыкии можно добраться либо на лошади, либо на верблюде. Докажите с помощью индукции, что можно выбрать один вид животного так, что используя только его по-прежнему можно добраться от любой стоянки до любой другой. Будем считать, что в Калмыкии неизвестное число животноводческих стоянок.